\chapter{Σχετικές εργασίες}
\label{ch:related}

\review[!]{\textbf{Το κεφάλαιο χρειάζεται τη μεγαλύτερη δουλειά από σένα.} Οι εργασίες βρέθηκαν με αναζήτηση και οι περιγραφές τους βασίζονται σε περιλήψεις, όχι σε ανάγνωση του πλήρους κειμένου. Πριν κρατήσεις οποιαδήποτε πρόταση:
(1) \textbf{διάβασε κάθε εργασία} (τουλάχιστον περίληψη, μέθοδο, αποτελέσματα) και διόρθωσε όσα γράφω γι' αυτήν·
(2) επιβεβαίωσε τα στοιχεία στο \code{references.bib} (συγγραφείς, τόμος, σελίδες, DOI)·
(3) \textbf{πρόσθεσε περισσότερες εργασίες}: μια διπλωματική θέλει συνήθως 25--40 αναφορές συνολικά. Λέξεις-κλειδιά στον \code{Οδηγός_Διπλωματικής.md} §5.}

Η εργασία βρίσκεται στην τομή τεσσάρων περιοχών: της μικτής πραγματικότητας για επιθεώρηση και παρακολούθηση κατασκευών, του BIM στο εργοτάξιο, της γεωαναφοράς σε επαυξημένη πραγματικότητα, και της οπτικοποίησης δεδομένων στη θέση τους (\en{situated visualization}). Το κεφάλαιο εξετάζει καθεμιά και κλείνει με μια σύγκριση.

\section{Μικτή πραγματικότητα για επιθεώρηση και παρακολούθηση κατασκευών}
\label{sec:rw-shm}

Η ιδέα να συνδεθεί η οπτική επιθεώρηση με ψηφιακά εργαλεία μέσα από συσκευές AR έχει διατυπωθεί ρητά στη βιβλιογραφία της παρακολούθησης κατασκευών. Οι \en{Mascareñas} κ.ά. προτείνουν ένα πλαίσιο λογισμικού για εφαρμογές επιθεώρησης υποδομών με AR. Ξεκινούν από το ότι η οπτική επιθεώρηση παραμένει η κύρια μέθοδος, και δείχνουν πώς μια συσκευή κεφαλής μπορεί να καταγράφει τρισδιάστατες μετρήσεις και να τις συνδέει με την κατασκευή \cite{mascarenas2021ar}.

Πιο κοντά στην παρούσα εργασία είναι συστήματα που φέρνουν \emph{δεδομένα αισθητήρων} στο πεδίο. Οι \en{Fawad} κ.ά. συνδέουν το σύστημα παρακολούθησης μιας νέας γέφυρας (αισθητήρες παραμόρφωσης, στάθμης, επιταχυνσιόμετρα, μετεωρολογικός σταθμός) με μια εφαρμογή AR σε μηχανή παιχνιδιών για το HoloLens. Οι μετρήσεις έρχονται από διαδικτυακή πλατφόρμα και η απεικόνιση γίνεται είτε επιτόπου είτε εξ αποστάσεως \cite{fawad2024bridge}. Οι \en{Awadallah} κ.ά. παρουσιάζουν σύστημα SHM με AR στο HoloLens~2 με έμφαση στην ακριβή ευθυγράμμιση του τρισδιάστατου μοντέλου με την κατασκευή, μέσω αναγνώρισης του μοντέλου από την κάμερα \cite{awadallah2025ar}.

Μια τρίτη ομάδα εργασιών εστιάζει στη διαχείριση της πληροφορίας επιθεώρησης και όχι στις συνεχείς μετρήσεις. Οι \en{Nguyen} κ.ά. αναπτύσσουν την εφαρμογή HoloBridge για το HoloLens. Βασίζεται σε μοντέλο πληροφοριών γέφυρας (BIM) και επιτρέπει την καταγραφή και ενημέρωση ευρημάτων επιθεώρησης και συντήρησης. Δοκιμάστηκε σε υπάρχουσα γέφυρα στη Νότια Κορέα \cite{nguyen2022holobridge}.

\section{BIM και μικτή πραγματικότητα στο εργοτάξιο}
\label{sec:rw-bim}

Η υπέρθεση μοντέλων σχεδιασμού στον πραγματικό χώρο είναι από τις πιο μελετημένες εφαρμογές της MR στις κατασκευές. Οι \en{Chalhoub} και \en{Ayer} μελέτησαν τη χρήση MR για την επικοινωνία του σχεδιασμού ηλεκτρολογικών εγκαταστάσεων σε εργάτες, και κατέγραψαν επιπτώσεις στην απόδοση της εργασίας σε σχέση με τα συμβατικά σχέδια \cite{chalhoub2018mixed}. Στην παρούσα εργασία το μοντέλο (BIM, CAD ή φωτογραμμετρία) δεν είναι ο κύριος σκοπός, αλλά το \emph{πλαίσιο} μέσα στο οποίο διαβάζονται οι μετρήσεις. Τοποθετείται με την ίδια ευθυγράμμιση με τους αισθητήρες, ώστε να μην χρειάζεται δεύτερη διαδικασία.

\todo{Πρόσθεσε 3--5 εργασίες για BIM σε AR/MR στο εργοτάξιο (π.χ. έλεγχος προόδου, σύγκριση as-built/as-designed) και για ψηφιακά δίδυμα υποδομών.}

\section{Γεωαναφορά σε επαυξημένη πραγματικότητα}
\label{sec:rw-georef}

Η τοποθέτηση γεωαναφερμένων γραφικών σε εξωτερικούς χώρους απασχολεί τη βιβλιογραφία εδώ και δύο δεκαετίες. Οι \en{Behzadan} και \en{Kamat} περιέγραψαν μέθοδο ευθυγράμμισης γραφικών κατασκευής σε κινητό σύστημα AR εξωτερικού χώρου, με θέση από GPS και προσανατολισμό από αισθητήρες κεφαλής \cite{behzadan2007georeferenced}. Οι \en{Schall} κ.ά. ανέπτυξαν σύστημα AR χειρός για την απεικόνιση υπόγειων δικτύων υποδομής από δεδομένα GIS, με εντοπισμό RTK-GPS \cite{schall2009handheld}. Πιο πρόσφατα, οι \en{Messi} κ.ά. συνδυάζουν πολλαπλούς μηχανισμούς ευθυγράμμισης χωρίς δείκτες (GNSS-RTK και όραση υπολογιστή) για εφαρμογές διαχείρισης εγκαταστάσεων σε εσωτερικούς και εξωτερικούς χώρους. Αξιολογούν ποσοτικά την ακρίβεια και δείχνουν ότι απαιτείται συχνή επαναευθυγράμμιση για ακρίβεια της τάξης του δεκατόμετρου \cite{messi2024seamless}.

Τα συστήματα αυτά βασίζονται σε GNSS, που δίνει ακρίβεια εκατοστών μόνο με εξοπλισμό RTK, δεν λειτουργεί σε κλειστούς χώρους, σήραγγες ή κάτω από καταστρώματα γεφυρών, και προσθέτει υλικό που πρέπει να συνδεθεί με τη συσκευή. Η παρούσα εργασία ακολουθεί την προσέγγιση των \emph{σημείων ελέγχου} της τοπογραφίας. Ο χρήστης ακουμπά με το χειριστήριο σημεία γνωστών συντεταγμένων, συνήθως τους ίδιους τους αισθητήρες, και το σύστημα λύνει τον μετασχηματισμό. Δεν χρειάζεται επιπλέον εξοπλισμός, λειτουργεί σε κλειστούς χώρους, και η ακρίβεια μπορεί να προβλεφθεί και να ελεγχθεί.

\section{Οπτικοποίηση δεδομένων στη θέση τους}
\label{sec:rw-situated}

Η απεικόνιση δεδομένων στον χώρο όπου αυτά έχουν νόημα μελετάται στην περιοχή της οπτικοποίησης ως \en{situated visualization} και \en{situated analytics}. Οι \en{Willett} κ.ά. διακρίνουν τις \emph{situated} αναπαραστάσεις, που απλώς βρίσκονται κοντά στο φυσικό αντικείμενο, από τις \emph{embedded}, που ενσωματώνονται σε αυτό \cite{willett2017embedded}. Οι \en{ElSayed} κ.ά. παρουσίασαν εργαλεία \en{situated analytics} με AR \cite{elsayed2016situated}. Δύο πρόσφατες επισκοπήσεις οργανώνουν το πεδίο: οι \en{Bressa} κ.ά. εξετάζουν πώς χρησιμοποιείται η έννοια της «θέσης» (χώρος, χρόνος, τόπος, δραστηριότητα, κοινότητα) \cite{bressa2022situated}, και οι \en{Shin} κ.ά. ταξινομούν συστήματα \en{situated analytics} σε τρεις διαστάσεις \cite{shin2023situated}. Ευρύτερα, το πεδίο του \en{immersive analytics} μελετά την ανάλυση δεδομένων σε εμβυθιστικά περιβάλλοντα \cite{marriott2018immersive}.

Με την ορολογία των \en{Willett} κ.ά., τα εικονίδια των αισθητήρων και οι ράβδοι ορίων της παρούσας εργασίας είναι \emph{situated}, ενώ ο χάρτης θερμότητας πάνω στο μοντέλο και τα διανύσματα μετακίνησης πάνω στους στόχους πλησιάζουν τις \emph{embedded} αναπαραστάσεις.

\review{Επιβεβαίωσε ότι η κατάταξη situated/embedded που κάνω εδώ συμφωνεί με τους ορισμούς του \cite{willett2017embedded}. Είναι καλό σημείο για συζήτηση στην εξέταση, άρα πρέπει να το καταλαβαίνεις.}

\section{Σύγκριση}
\label{sec:rw-comparison}

Ο \cref{tab:rw-comparison} συγκρίνει τις εργασίες που βρίσκονται πιο κοντά στο προτεινόμενο σύστημα.

\begin{table}[htbp]
  \centering
  \footnotesize
  \caption{Σύγκριση σχετικών εργασιών.}
  \label{tab:rw-comparison}
  \setlength{\tabcolsep}{4pt}
  \begin{tabularx}{\textwidth}{@{}>{\raggedright\arraybackslash}X
      >{\raggedright\arraybackslash}p{1.6cm} >{\raggedright\arraybackslash}p{1.9cm}
      >{\raggedright\arraybackslash}p{2.6cm} >{\raggedright\arraybackslash}p{1.7cm}
      >{\raggedright\arraybackslash}p{1.6cm}@{}}
    \toprule
    Εργασία & Συσκευή & Ζωντανά δεδομένα αισθητήρων & Ευθυγράμμιση & Μοντέλο 3Δ & Ποσοτική αξιολόγηση ακρίβειας \\
    \midrule
    \cite{fawad2024bridge} & HoloLens & Ναι (web) & \fillin{;} & BIM & \fillin{;} \\
    \cite{awadallah2025ar} & HoloLens 2 & \fillin{;} & Αναγνώριση μοντέλου από κάμερα & Ναι & Ναι \\
    \cite{nguyen2022holobridge} & HoloLens & Όχι (επιθεώρηση) & \fillin{;} & BIM & \fillin{;} \\
    \cite{schall2009handheld} & Χειρός & Όχι (GIS) & RTK-GPS & Δίκτυα & \fillin{;} \\
    \cite{messi2024seamless} & \fillin{;} & Όχι (FM) & GNSS-RTK + όραση & BIM & Ναι \\
    \midrule
    \textbf{Παρούσα εργασία} & Quest 3 & Ναι (REST, polling, push) & Σημεία ελέγχου + spatial anchors & OBJ / DXF / IFC & Ναι \\
    \bottomrule
  \end{tabularx}
\end{table}

\review{Συμπλήρωσε τα \fillin{;} \textbf{μόνο αφού διαβάσεις} τις εργασίες. Αν μια εργασία δεν καλύπτει μια στήλη, γράψε «Όχι» ή «—».}

Από τη σύγκριση προκύπτει ότι οι εργασίες που φέρνουν δεδομένα αισθητήρων στο πεδίο σπάνια αξιολογούν ποσοτικά την ακρίβεια τοποθέτησης. Αντίστροφα, οι εργασίες που εστιάζουν στην ευθυγράμμιση δεν αφορούν ζωντανά δεδομένα παρακολούθησης. Η παρούσα εργασία συνδυάζει: ζωντανά δεδομένα από πλατφόρμα παρακολούθησης, ευθυγράμμιση χωρίς επιπλέον εξοπλισμό με πρόβλεψη σφάλματος, χαμηλού κόστους αυτόνομη συσκευή, και ποσοτική αξιολόγηση στο πεδίο.

\review[!]{Η τελευταία παράγραφος είναι \textbf{ο ισχυρισμός πρωτοτυπίας} της εργασίας. Ισχύει μόνο αν τον στηρίζει η βιβλιογραφία που θα διαβάσεις. Αν βρεις εργασία που κάνει ήδη τον ίδιο συνδυασμό, διατύπωσέ τον πιο μετριοπαθώς.}
